\documentclass[11pt]{article}
\usepackage[a4paper,margin=0.8in]{geometry}
\usepackage{amsmath,amssymb,booktabs,array,graphicx}
\usepackage[colorlinks=true,linkcolor=blue,urlcolor=blue]{hyperref}
\setlength{\emergencystretch}{2em}
\title{MovingThreeFluidSim: a coupled banded fluid step}
\author{Mathematica derivation and C++ implementation}
\date{22 September 2026}
\begin{document}
\maketitle

\section{Scope and choices}
The new class is \texttt{MovingThreeFluidSim}. Its files are
\texttt{moving\_\allowbreak three\_\allowbreak fluid.hpp} and
\texttt{moving\_\allowbreak three\_\allowbreak fluid.cpp}, leaving \texttt{ThreeFluidSim}
unchanged. The starting code is upstream revision \texttt{49dd74e}.
The alternating-flux implementation is Code revision \texttt{af639ef}.
The previous tidal-force PR has been closed. This note derives one
linearly implicit, first-order time step for three isotropic fluids,
including self-gravity, an optional fixed outward acceleration $qr$,
conduction, intercomponent thermal exchange and binary heating.
Mass conversion by binary formation and an imposed stripping sink are
off in this first derivation; they are not silently counted as outflow.
The retained $c_1,c_4$ terms use the existing small-relative-streaming
thermal closure. Intercomponent mean-velocity drag and its compensating
heat are not derived here; appreciable relative bulk motion would require
that additional local physical closure.

The nonlinear dispersions use exactly the first-order expansions in
\texttt{conduction\_compiler.wl}. The conservative spatial fluxes are
not coefficient-identical to the old nonconservative log-grid stencil:
this deliberate change is needed to make boundary budgets telescope.
We also require an upwind treatment of moving fluid, absent from the
hydrostatic solver. The full state uses the requested primitive-variable
ordering. One banded solve determines a primitive predictor; a local
algebraic recovery preserves the conservative momentum/energy update.
There is no nonlinear iteration, timestep retry, or hydrostatic projection
inside the proposed step.

\paragraph{Geometry.}
For the first implementation use a fixed common radial mesh, which can
be logarithmic away from the centre. ``Moving'' denotes radial fluid
motion. The tidal radius moves diagnostically; it is not the last grid
face. A dynamically moving mesh would require geometry state or an
explicit mesh law in addition to the requested $12N$ fluid entries.
It is not hidden in those entries. Fixed geometry is the minimal closed
choice for this derivation and allows material to leave the outer shell.

\paragraph{Thermal boundary choice.}
``Freely flowing heat'' is not a unique boundary condition. We choose an
absorbing Robin condition on $\sqrt u$, specified in
Sec.~\ref{sec:outer}, rather than the insulating condition $u_r=0$.
Its extrapolation length is an explicit boundary-model choice, not
something derived from $q$. This is the approved baseline boundary model.

\section{Equations, units, and state}
Use fixed units $M_0=4\pi\rho_0r_0^3$, $u_0=4\pi G\rho_0r_0^2$,
$P_0=\rho_0u_0$, $L_0=M_0u_0/t_0$, and the existing relaxation-time
unit $t_0$. In all equations below,
\begin{equation}
 r=\frac{r_{\rm phys}}{r_0},\quad t=\frac{t_{\rm phys}}{t_0},\quad
 v=\frac{v_{\rm phys}t_0}{r_0},\quad
 \epsilon=\frac{r_0^2}{u_0t_0^2},\quad
 q=\frac{q_{\rm phys}}{4\pi G\rho_0}.
\end{equation}
Density, enclosed mass and specific energy are divided by
$\rho_0,M_0,u_0$ respectively. Keep $\gamma=5/3$ and
$\beta=\gamma-1=2/3$, matching the remote code; no new binary adiabatic
index is introduced. For each fluid $f\in\{s,b,d\}$,
\begin{align}
 P_f&=\beta\rho_fu_f,&
 M_f'&=r^2\rho_f,&
 g&=\frac{\sum_fM_f}{r^2}-qr,\\
 \partial_t\rho_f+\frac{1}{r^2}\partial_r(r^2\rho_fv_f)&=0,&
 \epsilon(\partial_t+v_f\partial_r)v_f&=-\frac{\partial_rP_f}{\rho_f}-g,
 \label{eq:pde}\\
 (\partial_t+v_f\partial_r)u_f+\frac{P_f}{\rho_f r^2}
 \partial_r(r^2v_f)&=-\frac{\partial_rL_f}{\rho_f r^2}+Q_f.
\end{align}
The physical one-dimensional dispersion satisfies
$\sigma_f^2/u_0=2u_f/3$. The sound speed in radius-per-code-time units is
$c_f=\sqrt{\gamma\beta u_f/\epsilon}$. Omitting $\epsilon$ in radial
momentum or bulk kinetic energy would change the time normalization.

Faces are $R_0=0<R_1<\cdots<R_N=B$, cell centres are
$r_i=(R_i+R_{i+1})/2>0$, and
\begin{equation}
 V_i=\frac{R_{i+1}^3-R_i^3}{3},\quad A_i=R_i^2,\quad
 \eta_i=\frac{r_i^3-R_i^3}{R_{i+1}^3-R_i^3},\quad
 a_i=\frac{A_{i+1}-A_i}{V_i}.
\end{equation}
Use cell-mean $\rho_{f,i}$ and outer-face enclosed mass $M_{f,i}=M_f(R_{i+1})$.
Thus $M$ and the other entries in one block are staggered geometrically,
though contiguous in memory. The exact discrete mass definition and
centre gravity approximation are
\begin{align}
 M_{f,i}-M_{f,i-1}&=V_i\rho_{f,i},\qquad M_{f,-1}=0,
 \label{eq:massconstraint}\\
 \overline M_i&=\eta_i\sum_fM_{f,i}+(1-\eta_i)\sum_fM_{f,i-1},&
 g_i&=\overline M_i/r_i^2-qr_i.
\end{align}
This mass interpolation is exact within a uniform-density cell.
The implemented, bandwidth-optimized vector in each cell is
\begin{equation}
 X_i=(u_s,v_s,\rho_s,u_b,v_b,\rho_b,u_d,v_d,\rho_d,M_s,M_b,M_d)^T.
 \label{eq:ordering}
\end{equation}
For $f=0,1,2$ and logical field $k=0,1,2,3\leftrightarrow(\rho,v,u,M)$,
the scalar index is $12i+3f+2-k$ for $k<3$, and $12i+9+f$ for $k=3$.
Equation rows use the same permutation: energy, momentum, continuity
for each species, then the three enclosed-mass constraints. The RHS
must be permuted with the rows. Formulas below retain logical
$(\rho,v,u)$ ordering within each displayed $3\times3$ block for readability;
their placement uses the index map above.
The redundant enclosed masses are constrained unknowns, not separate
mass reservoirs. Keeping Eq.~\eqref{eq:massconstraint} in the banded
system avoids a dense all-inner-cells gravitational Jacobian.

\section{Conservative variables and explicit Jacobians}
For $x=(\rho,v,u)^T$, define
\begin{equation}
 W(x)=\begin{pmatrix}\rho\\\rho v\\\rho(u+\epsilon v^2/2)\end{pmatrix},\qquad
 F(x)=\begin{pmatrix}\rho v\\\rho v^2+\beta\rho u/\epsilon\\
 \rho v(\gamma u+\epsilon v^2/2)\end{pmatrix}.
\end{equation}
The three rows are continuity, radial momentum, and random-plus-bulk
energy. Mathematica differentiates these to give
\begin{align}
 T=\frac{\partial W}{\partial x}
 &=\begin{pmatrix}1&0&0\\v&\rho&0\\
 u+\epsilon v^2/2&\epsilon\rho v&\rho\end{pmatrix},
 \label{eq:timejac}\\
 H=\frac{\partial F}{\partial x}
 &=\begin{pmatrix}
 v&\rho&0\\
 v^2+\beta u/\epsilon&2\rho v&\beta\rho/\epsilon\\
 v(\gamma u+\epsilon v^2/2)&\rho(\gamma u+3\epsilon v^2/2)&\gamma\rho v
 \end{pmatrix}.
 \label{eq:fluxjac}
\end{align}
All unstarred coefficients below are evaluated at time $n$.
With $\delta x=x^*-x^n$, use
$W^*=W^n+T\delta x$ and $F^*=F^n+H\delta x$.
This is a single Taylor-linearized backward step, not an exact solution
of nonlinear backward Euler.

\subsection{Alternating acoustic flux and hydrostatic balance}
The former centred flux had a checkerboard null mode: a cell-alternating
pressure is invisible to its centred gradient. Velocity viscosity alone
does not remove that defective coupling. Heating-off scans also excite
the mode during contraction; it is not a dilute-fluid or heating-specific
effect. We replace the centred acoustic traces, not the fluid equations.

Let $x_L=x_i$, $x_R=x_{i+1}$ be the \emph{raw} primitive states,
$\bar\rho=(\rho_L+\rho_R)/2$, $\bar u=(u_L+u_R)/2$, $V=v_L$,
and $p_R=\beta\rho_Ru_R$. Define the consistent alternating flux
\begin{equation}
 F^{\rm alt}(x_L,x_R)=
 \begin{pmatrix}
 \bar\rho V\\
 \bar\rho V^2+p_R/\epsilon\\
 \bar\rho V(\bar u+\epsilon V^2/2)+p_RV
 \end{pmatrix}.
 \label{eq:alternating}
\end{equation}
The pressure work uses the \emph{same} pressure and velocity traces.
The choice of left velocity and right pressure is fixed, not switched
according to the sign of the velocity. Advective stabilization is retained
separately. The flux equals $F(x)$ when both states equal $x$; its spatial
accuracy is first order. This is an alternating-trace finite-volume
scheme, not a change of the stored velocity to a staggered-grid variable.

Its primitive Jacobians, including every nonzero entry, are
\begin{align}
 H^{\rm alt}_L&=
 \begin{pmatrix}
 V/2&\bar\rho&0\\
 V^2/2&2\bar\rho V&0\\
 V(\bar u+\epsilon V^2/2)/2&
 \bar\rho(\bar u+3\epsilon V^2/2)+p_R&\bar\rho V/2
 \end{pmatrix},\\
 H^{\rm alt}_R&=
 \begin{pmatrix}
 V/2&0&0\\
 V^2/2+\beta u_R/\epsilon&0&\beta\rho_R/\epsilon\\
 V(\bar u+\epsilon V^2/2)/2+\beta u_RV&0&
 \bar\rho V/2+\beta\rho_RV
 \end{pmatrix}.
\end{align}
Freeze $\lambda_0=\max(|v_L|,|v_R|)$,
$\lambda_a=\max(|v_L|+c_L,|v_R|+c_R)$ and
$\kappa=(\lambda_a-\lambda_0)\bar\rho$. The numerical flux and its
complete affine expansion are
\begin{align}
 \widehat F^n&=F^{\rm alt}
 -\frac{\lambda_0}{2}(W_R-W_L)
 -\frac{\kappa}{2}
 \begin{pmatrix}0\\v_R-v_L\\\epsilon(v_R^2-v_L^2)/2\end{pmatrix},\\
 J_L&=H^{\rm alt}_L+\frac{\lambda_0}{2}T_L+\frac{\kappa}{2}K(v_L),&
 J_R&=H^{\rm alt}_R-\frac{\lambda_0}{2}T_R-\frac{\kappa}{2}K(v_R),\\
 \widehat F^*&=\widehat F^n+J_L\delta x_L+J_R\delta x_R,
 \label{eq:faceflux}
\end{align}
where $K(v)=\left(\begin{smallmatrix}0&0&0\\0&1&0\\0&\epsilon v&0\end{smallmatrix}\right)$.
The energy component includes the work of the velocity stress.
These $J_L,J_R$ replace the former centred-flux Jacobians in every block
and RHS below. The three-cell stencil, storage ordering, half-bandwidths
$(13,14)$ and $O(N)$ cost are unchanged.

\paragraph{Why the alternating mode is damped.}
For a uniform static background on a uniform periodic grid, without
gravity or heating, the linear acoustic equations become
\begin{align}
 \partial_t\delta\rho&=-\rho_0D_-\delta v,&
 \partial_t\delta p&=-\gamma p_0D_-\delta v,\\
 \partial_t\delta v&=-\frac{D_+\delta p}{\epsilon\rho_0}
                  +\nu\Delta_h\delta v,&
 \nu&=\frac{ch}{2},\qquad c^2=\frac{\gamma p_0}{\epsilon\rho_0}.
\end{align}
Here $D_-v_i=(v_i-v_{i-1})/h$ and $D_+p_i=(p_{i+1}-p_i)/h$.
They are negative adjoints and satisfy $D_-D_+=\Delta_h$, whose
symbol at the checkerboard wavenumber is $-4/h^2$, not zero.
The quadratic acoustic energy
$\sum_i h[\epsilon\rho_0(\delta v_i)^2/2+
(\delta p_i)^2/(2\gamma p_0)]$ is nonincreasing. Backward Euler adds
temporal damping. For the checkerboard mode, with $z=c\Delta t/h$,
the two amplification eigenvalues have modulus
\begin{equation}
 |G_\pi|=(1+2z+4z^2)^{-1/2}<1\qquad(z>0).
\end{equation}
Both acoustic variables therefore decay; the large-step limit is zero.
This is not a promise that a finite step annihilates every perturbation,
nor a nonlinear stability theorem for a stratified, self-gravitating,
open domain. Entropy perturbations and resolved physical waves must not
be indiscriminately filtered away.

\paragraph{Reference balance without reconstructed primitive states.}
No frozen density/energy offsets are applied to face states. They are
unnecessary for equilibrium: the advective dissipation vanishes when
$v=0$, and the velocity stress vanishes as well. Retaining those offsets
in an expanding exterior can corrupt its thermal transport.

For a supplied $q=0$ reference, use $P^e_{i+1/2}=P^e_{i+1}$ and
$P^e_{i-1/2}=P^e_i$ at interior faces, matching the new right-pressure
trace. Define
\begin{equation}
 C^e_{f,i}=\frac{A_{i+1}P^e_{f,i+1/2}-A_iP^e_{f,i-1/2}}{\epsilon V_i}
 -\frac{a_iP^e_{f,i}-\rho^e_{f,i}g_i^e}{\epsilon},
 \qquad a^e_{f,i}=\frac{\epsilon C^e_{f,i}}{\rho^e_{f,i}}.
 \label{eq:wb}
\end{equation}
The fixed acceleration $a^e$ enters momentum as $\rho_fa^e/\epsilon$
and energy as $\rho_fv_fa^e$. It vanishes as a force density when the
fluid empties. At the first face the area is zero. At the last face the
reference quadrature uses $P^e_{N-1}$, but the actual boundary remains
the vacuum flux below: a truncated positive-pressure sphere is not a
global equilibrium against vacuum. The interior reference is balanced
when thermal sources vanish. The physical $qr$ force is never absorbed
in $a^e$. With no reference, set $a^e=0$. The correction is numerical
quadrature, not an additional physical force, and must converge away.

\section{Gravity, conduction and the existing local sources}
\subsection{Dynamical cell source}
Apart from thermal sources, the finite-volume source is
\begin{equation}
 S_f^{\rm dyn}=\begin{pmatrix}
 0\\(a_i\beta\rho_fu_f-\rho_f\widetilde g_{f,i})/\epsilon\\-\rho_fv_f\widetilde g_{f,i}
 \end{pmatrix}.
\end{equation}
Here $\widetilde g_{f,i}=g_i-a^e_{f,i}$. The geometric pressure term is essential in conservative spherical
momentum. Its quadrature integrates constant pressure exactly.
The primitive and enclosed-mass derivatives are, explicitly,
\begin{align}
 D_f=\frac{\partial S_f^{\rm dyn}}{\partial x_f}
 &=\begin{pmatrix}0&0&0\\
 (a_i\beta u_f-\widetilde g_{f,i})/\epsilon&0&a_i\beta\rho_f/\epsilon\\
 -v_f\widetilde g_{f,i}&-\rho_f\widetilde g_{f,i}&0\end{pmatrix},\\
 G_f^{(0)}=\frac{\partial S_f^{\rm dyn}}{\partial M_{g,i}}
 &=\begin{pmatrix}0\\-\rho_f\eta_i/(\epsilon r_i^2)\\-\rho_fv_f\eta_i/r_i^2\end{pmatrix},&
 G_f^{(-)}=\frac{\partial S_f^{\rm dyn}}{\partial M_{g,i-1}}
 &=\begin{pmatrix}0\\-\rho_f(1-\eta_i)/(\epsilon r_i^2)\\-\rho_fv_f(1-\eta_i)/r_i^2\end{pmatrix}.
 \label{eq:sourcejac}
\end{align}
The two $G$ vectors apply to every species $g$. At the centre set
$G_f^{(-)}=0$. There is no admissibility condition $\overline M_i>qr_i^3$:
positive and negative net acceleration are both allowed.

\subsection{Dispersion expansion and conservative conduction}
Exactly as in the reference Mathematica calculation,
\begin{align}
 z_f(u_f^*)&=\sqrt{u_f^n}+\frac{\delta u_f}{2\sqrt{u_f^n}}
 =\frac{\sqrt{u_f^n}}2+\frac{u_f^*}{2\sqrt{u_f^n}},\\
 w_f(u_f^*)&=\frac1{\sqrt{u_f^n}}-
 \frac{\delta u_f}{2(u_f^n)^{3/2}}
 =\frac3{2\sqrt{u_f^n}}-\frac{u_f^*}{2(u_f^n)^{3/2}}.
 \label{eq:dispersion}
\end{align}
Freeze face conductivity at time $n$. With harmonic face density and
distance $d=r_{i+1}-r_i$,
\begin{equation}
 k_{f,i+1/2}=\frac{c_{2,f}A_{i+1}}{d}
      \frac{2\rho_{f,i}^n\rho_{f,i+1}^n}{\rho_{f,i}^n+\rho_{f,i+1}^n},\qquad
 L_{f,i+1/2}^*=-k_{f,i+1/2}(z_{f,i+1}-z_{f,i}).
 \label{eq:conductiveflux}
\end{equation}
This discretizes the same continuum flux
$L_f=-c_{2,f}r^2\rho_f\partial_r\sqrt{u_f}$ as the existing conduction
law. Adjacent cells use the same luminosity, with opposite signs.
All cells, including the first and last, retain their energy equations
and local sources. Neither boundary is a fixed-$u$ reservoir row.

\subsection{Exchange and binary heating}
Use the reference particle masses $m_f$ and define frozen
\begin{equation}
 a_{fg}=\frac{c_{1,fg}\rho_g^n}
 {m_s(u_f^n+u_g^n)^{3/2}},\quad g\ne f,\qquad
 b_{fg}=\rho_f^n c_{4,fg}\rho_g^n.
\end{equation}
The energy-density source, with the old density factors frozen just as
in the existing conduction step, is
\begin{equation}
 Q_f^{E,*}=-\rho_f^n\sum_{g\ne f}a_{fg}(m_fu_f^*-m_gu_g^*)
       +\sum_{(g,d)\in\mathcal H_f}b_{fg}w_d(u_d^*),
\end{equation}
where
$\mathcal H_s=\{(b,s)\}$,
$\mathcal H_b=\{(s,s),(b,b),(d,d)\}$, and
$\mathcal H_d=\{(b,d)\}$.
Writing $Q_f^{E,*}=Q_f^{E,n}+\sum_g B_{fg}\delta u_g$, every local
temperature coefficient is
\begin{equation}
 \boxed{B_{fg}=-\delta_{fg}\rho_f^nm_f\sum_{h\ne f}a_{fh}
 +(1-\delta_{fg})\rho_f^na_{fg}m_g
 -\sum_{(h,d)\in\mathcal H_f}\frac{b_{fh}}{2(u_d^n)^{3/2}}\delta_{dg}.}
 \label{eq:thermaljac}
\end{equation}
For $g=f$ the second term is defined as zero. The old source is the
same expression before expansion, using $w_d=1/\sqrt{u_d^n}$.
In an absolute-variable solve the heating constant is
$\sum_{(g,d)\in\mathcal H_f}3b_{fg}/(2\sqrt{u_d^n})$, reproducing the
reference factor $3/2$. Symmetric $c_{1,fg}$ gives pairwise exchange
energy cancellation. With fixed density, no motion and no gravity work,
the local $c_1,c_4$ matrix reduces to the supplied conduction derivation
after dividing the energy row by $V_i\rho_{f,i}^n$.

\section{Centre and open outer boundary}
\subsection{Regular centre, not four replacement PDE rows}
At the physical face $R_0=0$, use even $\rho,u$ and odd $v$ under
reflection. This imposes $\rho_r=0$, $u_r=0$ and $v=0$ at the centre.
The first \emph{cell-centre} velocity need not vanish: a regular flow
has $v(r)=v_1r+O(r^3)$. Set central mass, momentum-area and energy-area
fluxes to zero and $L_f(0)=0$, and keep all first-cell evolution equations.
Replacing them by $\rho_0=\rho_1$, $u_0=u_1$ would discard a finite
cell's mass/energy balance.

Mass regularity gives
\begin{equation}
 M_f(r)=\frac{\rho_{f,c}}3r^3+O(r^5),\qquad
 rM_f'(r)-3M_f(r)=O(r^5),\qquad M_f(0)=0.
 \label{eq:robin}
\end{equation}
The limiting Robin relation is degenerate at exactly $r=0$; it is not
an additional independent boundary datum. In the adopted finite-volume
layout its first-cell counterpart is
\begin{equation}
 \boxed{M_{f,0}-V_0\rho_{f,0}=0,\quad V_0=R_1^3/3.}
\end{equation}
This supplies the requested regular mass closure without $0/0$ or
overconstraining the first-order enclosed-mass equation. There is no
additional mass boundary condition at $B$.

\subsection{Material outflow into vacuum}\label{sec:outer}
Use the ideal-gas Riemann problem between the last populated interior
state and vacuum, rather than prescribing $v$, $\rho$ and $P$ separately.
At the old interior state define
\begin{equation}
 c=\sqrt{\gamma\beta u/\epsilon},\qquad
 v_{\rm head}=v-c,\quad v_{\rm front}=v+\frac{2c}{\gamma-1}.
\end{equation}
Freeze the selected regime for one linearized step:
\begin{equation}
 F_B(x)=\begin{cases}
 F(x),&v_{\rm head}\ge0,\\
 F\!\left(\rho\chi^{2/(\gamma-1)},c\chi,u\chi^2\right),
 &v_{\rm head}<0<v_{\rm front},\\
 0,&v_{\rm front}\le0,
 \end{cases}\qquad
 \chi=\frac{2c+(\gamma-1)v}{(\gamma+1)c}.
 \label{eq:vacuum}
\end{equation}
The intermediate branch is the rarefaction sampled at the stationary
face.\footnote{For the Euler rarefaction invariants, see the
\href{https://www.clawpack.org/riemann_book/html/Euler.html}{Clawpack Riemann book, Euler chapter}.}
Its mass flux is nonnegative. It lets a pressure-supported boundary
expand even if its old mean velocity is zero. Gravity acts through the
cell source, not through an extra hydrostatic boundary row.

For the fan branch at $\gamma=5/3$, let
$R=\rho\chi^3$, $V=c\chi=(3c+v)/4$, $U=u\chi^2$.
The complete derivative needed by the banded matrix is $H_B=H(R,V,U)K$,
where
\begin{equation}
 K=\frac{\partial(R,V,U)}{\partial(\rho,v,u)}=
 \begin{pmatrix}
 \chi^3&3\rho\chi^2/(4c)&-3\rho\chi^2v/(8cu)\\
 0&1/4&3c/(8u)\\
 0&u\chi/(2c)&\chi^2-\chi v/(4c)
 \end{pmatrix}.
 \label{eq:vacjac}
\end{equation}
For the first and last branches, $H_B=H$ and $H_B=0$ respectively.
Thus $F_B^*=F_B^n+H_B\delta x_{N-1}$ couples only the last cell.
The compiler differentiates Eq.~\eqref{eq:vacuum} directly and checks
its result independently. At regime transitions or a vacuum cell,
use the appropriate one-sided/vacuum treatment, not a derivative through
$u=0$. A positive-state linearization is not itself a complete vacuum-cell
algorithm; that extension must be tested before evolving actual empty cells.

\subsection{Nonzero outward heat flux}
An insulating $u_r(B)=0$ would contradict the requested thermal outflow.
Choose instead
\begin{equation}
 \partial_r z(B)+\frac{z(B)}{\ell_{\rm th}}=0,\qquad
 z=\sqrt u,\quad\ell_{\rm th}>0.
\end{equation}
With $d_B=B-r_{N-1}$, linear interpolation to the face gives
\begin{equation}
 L_{f,B}^*=h_f z_{f,N-1},\qquad
 h_f=\frac{c_{2,f}B^2\rho_{f,N-1}^n}{\ell_{\rm th}+d_B}\ge0.
 \label{eq:thermalboundary}
\end{equation}
For this proposed baseline choose $\ell_{\rm th}=B$, keeping the
coefficient finite under refinement; retain $h_f$ symbolically in the
compiler for auditing alternatives. This is an absorbing exterior thermal
closure, not a derivation of a unique collisionless escape luminosity.
It allows energy loss without imposing a fixed last-cell temperature.
It is independent of advected enthalpy; both contributions enter the
energy ledger. Setting $h_f=0$ recovers a closed thermal test.

\section{Every block entry and the RHS}\label{sec:matrix}
Let $e_3=(0,0,1)^T$, $e_u=(0,0,1)^T$, and consider cell $i$.
The three fluid rows are
\begin{equation}
 V_iT_{f,i}\delta x_{f,i}+\Delta t\left[
 A_{i+1}\widehat F_{f,+}^*-A_i\widehat F_{f,-}^*
 +e_3(L_{f,+}^*-L_{f,-}^*)-V_iS_f^*\right]=0,
 \label{eq:assembledrow}
\end{equation}
where $S_f^*=S_f^{\rm dyn,*}+e_3Q_f^{E,*}$.
Each cell has a $12\times12$ diagonal block and two neighbouring blocks.
For compactness below, $A_-=A_i$, $A_+=A_{i+1}$,
$k_-=k_{f,i-1/2}$ and $k_+=k_{f,i+1/2}$.

The following $3\times3$ matrices specify \emph{every} primitive-column
entry in the first three rows of the species blocks; $T,H,D,B,G$ and
the interface Jacobians were explicitly listed above:
\begin{align}
 A^-_{fg}\big|_{x}
 &=\delta_{fg}\left[-\Delta t A_-J_{f,-,L}
 -\Delta t\frac{k_-}{2\sqrt{u_{f,i-1}^n}}e_3e_u^T\right],
 \label{eq:blockminus}\\
 A^0_{fg}\big|_{x}
 &=\delta_{fg}\left[V_iT_f+\Delta t(A_+J_{f,+,L}-A_-J_{f,-,R})
 -\Delta t V_iD_f
 +\Delta t\frac{k_-+k_+}{2\sqrt{u_{f,i}^n}}e_3e_u^T\right]
 -\Delta t V_i B_{fg}e_3e_u^T,\\
 A^+_{fg}\big|_{x}
 &=\delta_{fg}\left[\Delta t A_+J_{f,+,R}
 -\Delta t\frac{k_+}{2\sqrt{u_{f,i+1}^n}}e_3e_u^T\right].
 \label{eq:blockplus}
\end{align}
The enclosed-mass columns in those same three rows are
\begin{equation}
 A^-_{fg}\big|_M=-\Delta t V_iG_f^{(-)},\qquad
 A^0_{fg}\big|_M=-\Delta t V_iG_f^{(0)},\qquad
 A^+_{fg}\big|_M=0.
\end{equation}
The fourth (mass-constraint) row has only
\begin{equation}
 A^0_{(f,M),(f,\rho)}=-V_i,\qquad
 A^0_{(f,M),(f,M)}=1,\qquad
 A^-_{(f,M),(f,M)}=-1.
 \label{eq:massrow}
\end{equation}
All other entries in that row, and all entries not specified by these
formulas, are zero. This also states all cross-species couplings.

The increment RHS in the first three rows is
\begin{equation}
 \boxed{b_{f,i}=\Delta t\left[V_i(S_f^{{\rm dyn},n}+e_3Q_f^{E,n})
 -A_+\widehat F_{f,+}^n+A_-\widehat F_{f,-}^n
 -e_3(L_{f,+}^n-L_{f,-}^n)\right].}
 \label{eq:rhs}
\end{equation}
The fourth RHS is
$b_{f,M}=-(M_{f,i}^n-M_{f,i-1}^n-V_i\rho_{f,i}^n)$; it vanishes
for a consistent old state, but should not be silently discarded.

\paragraph{First cell.}
Delete all incoming-flux terms, set $k_-=0$, $G^{(-)}=0$,
$M_{f,-1}=0$, and delete the previous-mass coefficient in
Eq.~\eqref{eq:massrow}. Keep the full first-cell time and source terms.

\paragraph{Last cell.}
Replace $\widehat F_+^n$ by $F_B^n$, $J_{+,L}$ by $H_B$,
and set $J_{+,R}=0$. Replace the outgoing conduction diagonal $k_+$
by $h_f$, delete its next-cell term, and use
$L_+^n=h_f\sqrt{u_{f,N-1}^n}$ in the RHS. The inner-face terms and
mass-constraint row are unchanged. No next-cell unknown survives.
Different components may use different branches of Eq.~\eqref{eq:vacuum}.

The result is $A\delta X=b$. If one prefers to solve for the absolute
primitive predictor, the identical equation is
\begin{equation}
 AX^*=b_{\rm abs},\qquad b_{\rm abs}=b+AX^n.
\end{equation}
The companion exports both RHS conventions so that a missing constant
from a square-root expansion cannot be hidden by notation.

\section{Conservative recovery, outflow, and computational cost}
\subsection{One solve followed by local algebra}
Since $W$ is nonlinear in $(\rho,v,u)$, storing $x^n+\delta x$ directly
would not preserve the conservative time update exactly. Recover
\begin{align}
 \widetilde W_f&=W_f^n+T_f\delta x_f,&
 \rho_f^{n+1}&=\widetilde W_{f,1},\\
 v_f^{n+1}&=\widetilde W_{f,2}/\rho_f^{n+1},&
 u_f^{n+1}&=\widetilde W_{f,3}/\rho_f^{n+1}
                  -\epsilon(v_f^{n+1})^2/2.
 \label{eq:recovery}
\end{align}
Keep $M^{n+1}=M^n+\delta M$. Density is unchanged by recovery, so the
mass constraints remain satisfied. The recovery costs $O(1)$ per cell,
requires no matrix solve, and differs from the primitive predictor by
$O(|\delta X|^2)$. It is part of the proposed algorithm, not an optional
post hoc conservation repair. If final primitive values must be the
literal solution of a linear equation, exact nonlinear conservative
storage and a single solve cannot both be promised.

Use the \emph{same solved affine boundary flux}, before primitive
recovery, for the per-step outflow ledger:
\begin{align}
 \Delta M_{f,\rm out}&=\Delta t B^2\widehat F_{f,B,1}^*,\\
 \Delta E_{f,\rm adv}&=\Delta t B^2\widehat F_{f,B,3}^*,&
 \Delta E_{f,\rm heat}&=\Delta t L_{f,B}^*.
\end{align}
For zero central flux and no mass conversion,
\begin{equation}
 \boxed{M_{f,N-1}^{n+1}-M_{f,N-1}^n=-\Delta M_{f,\rm out}.}
\end{equation}
Do not recompute a nonlinear face flux from recovered snapshots and use
that different number in the budget. Summed internal conductive and
advective fluxes cancel. This guarantees the discrete \emph{mass} ledger,
not exact total gravitational energy conservation. Self-gravity work,
the external potential $-qr^2/2$, physical binary heating/binding budgets,
and gravitational energy carried by exported mass need their own
consistent energy audit.

Specific entropy is materially constant for smooth adiabatic flow in
the continuum, but this first-order dissipative discretization does not
preserve $P/\rho^{5/3}$ exactly. Heating/conduction and shocks change
entropy physically; numerical entropy must decrease under refinement.
This differs from the algebraic entropy-preserving hydrostatic projection.

\subsection{Bandwidth and the precise linear-complexity claim}
With Eq.~\eqref{eq:ordering}, symbolic enumeration gives
\begin{equation}
 n_{\rm eq}=12N,\qquad k_l=13,\qquad k_u=14.
\end{equation}
The original $(\rho_s,v_s,u_s,M_s,\ldots)$ layout gave $(k_l,k_u)=(19,13)$.
Moving masses to the block end and reversing the primitive triples gives
$(13,14)$. This is optimal among matching within-cell row/column
permutations: bidirectional neighbouring primitive couplings force
both half-bands to be at least 13, while three mutually coupled primitive
variables span at least two slots, forcing one half-band to be at least 14.
All centre and outer regimes are included in the symbolic sparsity audit.
A general pivoted band solve, not a dense solve and not an SPD solver,
has work $O(n_{\rm eq}k_l(k_l+k_u))=O(N)$ and memory $O(N)$.
LAPACK general-band storage needs $2k_l+k_u+1=41$ rows (previously 52), including pivot
fill-in.\footnote{Storage convention:
\href{https://netlib.org/lapack/explore-html/db/df8/group__gbsv_gaff55317eb3aed2278a85919a488fec07.html}{LAPACK DGBSV documentation}.}
Block assembly, local recovery, all budget sums and a scan for
$M(r)=qr^3$ are also $O(N)$. Do not eliminate the enclosed-mass constraints
into a dense Jacobian or compute each enclosed mass by a fresh inner sum.
The claim is \emph{per timestep}; refining the grid can require more
timesteps, so the full run need not scale linearly with $N$.

\subsection{Admissibility and later implementation}
One Taylor-linearized solve is not unconditionally positivity preserving.
Require positive density and internal energy, a valid populated-cell
temperature for all expansions, and nonnegative outward boundary mass
flux for the chosen vacuum model. Bound the step using wave propagation,
source changes and outgoing cell mass, as well as the existing fractional
$u$ accuracy criterion. If a step is invalid, report failure; do not
silently clip densities, count negative outflow as escape, or retry without
an agreed policy. An implicit solve does not guarantee accuracy at large
$\Delta t$ or cure low-Mach numerical diffusion. Vacuum activation,
reference source quadrature, sonic-branch changes and large-time-step
behaviour require numerical tests before production use.

The implemented class owns one contiguous fluid vector, fixed grid
metadata, band workspace and budget state. Evolving scalar ledgers can
be stored in a documented tail of the state vector (outside the $12N$
solve); geometry and physical conversion metadata must not be confused
with solved fluid entries. Keep evolution, initialization and observer
settings separate. Use named helpers for assembly, boundary fluxes,
recovery and validation, called through the shared-style \texttt{evolve()}
driver. Do not call hydrostatic relaxation after the moving step or
delete material at the diagnostic tidal radius. The implementation lives in
\texttt{moving\_three\_fluid.hpp/cpp}; initializer and comparison observer
settings remain separate.

\section{Companion Mathematica files and checks}
\texttt{moving\_three\_fluid\_compiler.wl} builds the nonlinear residuals,
their affine discretization, every coefficient by \texttt{Coefficient},
and the RHS by subtracting the coefficient--variable dot product, following
the method of the supplied conduction compiler. The supplied reference
file contains its expressions inside comments; the companion here is
executable and does not rely on uncommenting or changing that file.

Open \texttt{moving\_three\_fluid\_compiler.nb} and evaluate its first
cell. The resulting association has the cases
\texttt{centre}, \texttt{bulk}, \texttt{outer\_fan},
\texttt{outer\_supersonic}, and \texttt{outer\_vacuum}.
Each contains \texttt{Aminus}, \texttt{Adiag}, \texttt{Aplus},
\texttt{IncrementRHS} and \texttt{AbsoluteRHS}. The local column list
orders offsets $-1,0,+1$, then species, then $\rho,v,u,M$.
\texttt{moving\_three\_fluid\_entries.txt} lists every nonzero scalar
entry and each increment RHS; \texttt{moving\_three\_fluid\_coefficients.wl}
is a reloadable symbolic export. Unlisted entries are zero.

Checks include affine reconstruction of all rows, independent storage
and flux Jacobians, finite differences of the \emph{nonlinear} residuals
for every cell/boundary case, the scalar bandwidth, centre mass regularity,
telescoping of the assembled continuity equations, pairwise exchange
energy cancellation, and conservative primitive recovery.
The block formulas in Sec.~\ref{sec:matrix} are assembled independently
and compared entry by entry with the extracted coefficients.
The companion additionally exports \texttt{storageRowsByRegion}, containing
all blocks, variable lists and both RHS conventions in the optimized order.
The optimized half-bandwidths and affine identities are checked independently.
These are algebraic checks, not a completed hydrodynamic stability test.
The C++ implementation and sample comparisons are described in the companion
note \emph{MovingThreeFluidSim: implementation checks and sample comparisons}.
The earlier results note records the centred-flux checkpoint; the
updated oscillation checks below supersede its instability status.

\section{Contraction tests of the alternating flux}
The heating-off scan shows that strong contraction, not heating itself
or a small mass fraction, excites the old checkerboard mode. The new scheme
retains conduction, the timestep controller, gravity, outflow and source
coefficients. No retry, state floor, velocity filter or removal of trace
fluids is used. Two exploratory pressure-correction schemes were rejected:
an acoustic-strength face balance was inconsistent with the cell balance
and stalled; a weaker pressure correction barely changed the oscillation.
They are not retained as selectable implementations.

The alternating traces remove the acoustic null space. Removing the old
primitive offsets is also necessary: retaining them in this flux drove the
outermost nonconducting binary cell towards zero internal energy and stalled
the single-fluid run near $t=4.22$. With raw primitives and the corresponding
reference-source quadrature, all runs in Table~\ref{tab:oscillation} finish.
The physical outer material and heat boundary formulas are unchanged.

\begin{table}[ht]
\centering
\begin{tabular}{lrrrr}
\toprule
Case & $N$ & Target $t$ & Steps & Final stellar estimator\\
\midrule
Single / equal split &150&5&5537&$1.92\times10^{-4}$\\
AB-like &150&2.8&3206&$1.09\times10^{-4}$\\
Canonical &500&5&7237&$1.33\times10^{-5}$\\
Canonical refined &1000&5&7317&$3.05\times10^{-6}$\\
Concentrated DM, heating off &150&5&5967&$1.93\times10^{-4}$\\
Concentrated DM, heating off &300&5&6213&$3.05\times10^{-5}$\\
\bottomrule
\end{tabular}
\caption{The estimator is the median of
$|(2v_i-v_{i-1}-v_{i+1})/4|$ over $0.1<r_i<1$, in $r_0/t_0$ units.
It includes ordinary smooth curvature and is not a fixed-wavenumber norm.
Final-step overshoot is allowed.}\label{tab:oscillation}
\end{table}
The old single-fluid value was $0.119$, with repeated inward/outward
alternation; the new profile has one smooth transition from contraction
to expansion. The concentrated-DM setup has $\rho_{s,0}=1$,
$\xi_b=10^{-10}$, $\xi_d=1$, $\zeta_b=1$, $\zeta_d=0.3$,
$m_s=10^{-6}$, $m_b/m_s=2$, $m_d/m_s=10^{-10}$ and $c_1=c_4=0$.
Its old 150/300-cell estimators were approximately $0.53/1.26$.

\begin{figure}[ht]
\centering
\IfFileExists{../output/moving_oscillation_fix/single/velocity_oscillations.pdf}{%
\includegraphics[width=.9\textwidth]{../output/moving_oscillation_fix/single/velocity_oscillations.pdf}}{%
\fbox{\parbox{.85\textwidth}{Run the documented comparison and plotting commands
to generate \texttt{output/moving\_oscillation\_fix/single/velocity\_oscillations.pdf}.}}}
\caption{Old centred flux (red) and new alternating flux (blue) in the
single-fluid setup, including both trace components. Left: signed cell
velocities at common final time. Right: the second-difference estimator.
The positive residual estimator for a smooth curve is not evidence of
remaining checkerboarding. No time or amplitude rescaling is applied.}
\end{figure}

The C++ pressure/velocity perturbation tests, with conduction off and
$\epsilon=0.03$ or $10^{-9}$, reduce the respective alternating amplitudes
below $0.061$ of their initial values in one step with $c\Delta t/h=10$.
A separate no-heating/no-conduction test reaches $t=0.12$.
All 37 Mathematica checks pass. Matrix/RHS comparisons with C++ give scaled
errors $4.17\times10^{-16}$ and $1.68\times10^{-16}$. The unchanged mass
and gas-energy/work ledgers still close near roundoff.

These checks remove the observed instability; they do not establish a
fully converged collapse solution. In the concentrated-DM case the final
central densities are $123.4$ and $178.8$ at 150 and 300 cells. The
1000-cell canonical stellar density and dispersion remain within 10\%
of the hydrostatic reference throughout the tested interval, but the
coarse unequal-mass comparison still differs substantially. Before
production use, require spatial/time convergence and broader wave,
vacuum and tidal-escape tests. In particular, the raw ratio $v/\sqrt{2u/3}$
mixes velocity units: the ratio to physical dispersion is
$\sqrt{\epsilon}v/\sqrt{2u/3}$. Large raw plotted velocities alone
are not evidence of supersonic motion.
\end{document}
